\documentclass[../main]{subfiles}
\begin{document}
\ifSubfilesClassLoaded{\setcounter{page}{62}}{}
\section{Заключение}
\label{sec:11_zaklyuchenie}

Только теперь, когда мы изложили в общих чертах основные моменты коммунизма и показали те ростки, из которых он может вырасти как общественный строй, мы считаем себя вправе перейти к дефинициям. В противном случае, они были бы мертвыми, лишенными смысла абстракциями, а не ответом на вынесенный в название вопрос «что такое коммунизм?».

Коммунизм (коммунистическое общество) - это свободная ассоциация людей, разумно творящих свою жизнь, сознательно полагающих свою собственную свободную деятельность.

Коммунизм - это непосредственно-коллективная форма свободной деятельности, как форма общества, творящего своё собственное развитие, ставшая тотальной.

Коммунизм - это, как говорили Маркс и Энгельс в «Манифесте», «ассоциация, в которой свободное развитие каждого является условием свободного развития всех».

Пришло время опять повторить нисколько не устаревший лозунг:

ПРОЛЕТАРИИ ВСЕХ СТРАН, СОЕДИНЯЙТЕСЬ!
\end{document}
